\documentclass[12pt,english]{article}
\usepackage[T1]{fontenc}
\usepackage[latin9]{inputenc}
\usepackage[toc,page]{appendix}
\usepackage{geometry}
\geometry{verbose,tmargin=1.5in,bmargin=1.5in,lmargin=1in,rmargin=1in}
\usepackage{bm}
\usepackage{amstext}
\usepackage{comment}
\usepackage{amsmath}
\usepackage{hyperref}
\hypersetup{
     colorlinks = true,
     citecolor = navy,
     linkcolor = navy,
     urlcolor=navy
}
\usepackage[space]{grffile}
\usepackage{setspace}
\usepackage[normalem]{ulem}
\usepackage{xcolor}
\usepackage{soul}
\usepackage{dsfont}
\usepackage{refstyle}
\usepackage[capitalize,nameinlink]{cleveref}
\usepackage{setspace}
\usepackage{bbm}
\usepackage[tableposition=t]{caption} % ensures caption of longtables is normalsized.
\usepackage{mathpazo} % add Palatino font
\usepackage{longtable}
\usepackage{amssymb}
\usepackage{amsfonts}
\usepackage{amsmath}
%\usepackage[margin=1in]{geometry} % to set margins
\usepackage{scalefnt}
\usepackage{wrapfig}
\usepackage{graphicx}
%\usepackage{subfigure}
\usepackage{lscape}
\usepackage{pdflscape}
\usepackage{multirow}
\usepackage{array}
\usepackage{rotating}
\usepackage{bibunits}
\usepackage{fix-cm}
\usepackage{float}
\floatplacement{table}{H}
\usepackage{epsfig}
\usepackage{epstopdf}
\usepackage{refstyle}
\usepackage{booktabs}
\usepackage{tabularx}
\usepackage{array}

\newcolumntype{x}[1]{>{\raggedleft\arraybackslash}p{#1}}%
\newcolumntype{d}[1]{>{\centering\arraybackslash}p{#1}}%
\newcolumntype{y}[1]{>{\raggedright\arraybackslash}p{#1}}%
\usepackage{verbatim}
%\epstopdfsetup{suffix=}
\setcounter{MaxMatrixCols}{10}
\usepackage[longnamesfirst,authoryear]{natbib}
\newenvironment{proof}[1][Proof]{\noindent \textbf{#1.}}{\rule{0.5em}{0.5em}}
\setlength{\topmargin}{-0.80in} \setlength{\textheight}{9.50in}
\setlength{\evensidemargin}{.0in} \setlength{\oddsidemargin}{.0in}
\setlength{\textwidth}{6.85in}
\setlength{\voffset}{0.4cm}
\setlength{\hoffset}{-0.5cm}
\usepackage{setspace}
\usepackage{afterpage}
\usepackage[flushleft]{threeparttable}
\usepackage{tabulary}
%\usepackage[titletoc,title]{appendix}
\usepackage[bottom]{footmisc}
\usepackage{color}
	\definecolor{maroon}{RGB}{128,0,0}
	\definecolor{navy}{RGB}{0,0,165}
	\definecolor{lb}{RGB}{150,165,190}
	\definecolor{MyWhite}{RGB}{255, 255, 255}
\usepackage{colortbl}

\usepackage{caption}
\usepackage{subcaption}
\usepackage{placeins}
\usepackage{nicefrac}
%\usepackage{chngcntr}

\newcommand{\bit}{\begin{itemize}}
\newcommand{\eit}{\end{itemize}}
\usepackage{listings}
\def\cN{\mathcal{N}}
\newtheorem{theorem}{Theorem}
\newtheorem{remark}[theorem]{Remark}
\usepackage{mdframed}
\newmdtheoremenv{assumption}{Assumption}




\begin{document}


%\onehalfspacing


\title{Running the VAR for ShapirIO}
%Beggar or Prosper-thy-Neighbour? \\
%Fiscal Spillovers Across Firms and Countries
\maketitle

\paragraph{Purpose:} The purpose of this note is to introduce some details to run a VAR.

\section{What is a VAR}

A VAR is simply OLS with a vector of LHS variables. The ``V'' stands for Vector (meaning multiple LHS variables), ``A'' for Auto which simply means over time (so the RHS variables are going to be lags of the LHS variables).

So a 2-variable VAR will be like this:

\[
	\underbrace{\begin{bmatrix}
		\log{price}_{t}\\
		\log{quantity}_{t}
	\end{bmatrix}}_{\equiv \underbar{Y}_t}
	=
	\underbrace{\begin{bmatrix}
		a_{price,price} & q_{price quantity}\\
		a_{quantity,price}&a_{quantyt,quantity}
	\end{bmatrix}}_{\equiv A \text{ Parameter matrix}}
	\begin{bmatrix}
		\log{price}_{t-1}\\
		\log{quantity}_{t-1}
	\end{bmatrix}
	+\begin{bmatrix}
	e_{price,t}\\
	e_{quantity,t}
	\end{bmatrix}
\]

The parameter matrix $A$ can be estimated using the same formula for OLS adjusted for the fact that the LHS is a vector. We can stack $\underline{Y}_{t-1}$ into the following matrices:

\[
	\mathbf{Y}_{RHS}\equiv \begin{bmatrix}
		\underline{Y}_{0}\\
		\vdots\\
		\underline{Y}_{T-1}
	\end{bmatrix},\;\;\;\;\;\mathbf{Y}_{LHS}\equiv\begin{bmatrix}
		\underline{Y}_{1}\\
		\vdots\\
		\underline{Y}_{T}
	\end{bmatrix}
\]

\begin{align*}
	\hat{A}&=\left(\sum_t \underline{Y}_{t-1}'\underline{Y}_{t-1}\right)^{-1}\left(\sum_t\underline{Y}_{t-1}'\underline{Y}_t\right)&\text{Summation notation}\\
	&=\left(\mathbf{Y}_{RHS}'\mathbf{Y}_{RHS}\right)^{-1}\mathbf{Y}_{RHS}'\mathbf{Y}_{LHS}&\text{Matrix Notation}
\end{align*}

Note that Python has a package to do all this VAR stuff so I would not worry too much about the details of implementation.

Here is what I would like from you. 

\begin{enumerate}
	\item Using the price and quantity data we collected I want you to run VARs for every sector. This should be done with the BEA dataframe (not the one merged with the IO yet)
	\item The VAR above has a single lag on the RHS but one can run VARs with multiple lags. Since this is monthly data, I want you to run 24 lags to start with (we can consider lag length selection later).
	\item Then I want you to save the residuals from each sectoral regression in their own data frame.
	\item Then in that residual dataframe, I want you calculate an expenditure residual as the product of the price and quantity residuals.
\end{enumerate}

Here is some code that might help:
	
\begin{lstlisting}[label=some-code,caption=Some Code]
import pandas as pd
import numpy as np
import matplotlib.pyplot as plt


\# Import Statsmodels
from statsmodels.tsa.api import VAR
from statsmodels.tsa.stattools import adfuller
from statsmodels.tools.eval_measures import rmse, aic

model = VAR(df)
for i in [1,2,3,4,5,6,7,8,9]:
    result = model.fit(i)
    print('Lag Order =', i)
    print('AIC : ', result.aic)
    print('BIC : ', result.bic)
    print('FPE : ', result.fpe)
    print('HQIC: ', result.hqic, '\n')
\end{lstlisting}


\end{document}